\chapter{计算机辅助工艺设计（CAPP）}
\label{ch:capp}
\sourcepage{CAPP补充讲义}{3--5}

\section{为什么CAD与CAM之间还需要工艺设计}
CAD回答“零件应是什么形状、满足什么要求”；制造还必须回答“用什么毛坯、在哪台机床、怎样定位装夹、按什么顺序加工、用什么刀具和参数”。这些决定构成工艺设计。即使两个零件几何形状相同，材料、精度、批量或可用设备不同，也可能需要不同工艺路线。

\dfn{工艺规程与CAPP}{工艺规程是规定产品或零件制造过程和操作方法的技术文件，可表现为工艺过程卡、工序卡、工步卡及工序图等。CAPP（Computer-Aided Process Planning）利用计算机辅助工艺人员编制、管理和优化工艺规程。其核心是\textbf{把产品要求与制造资源转化为可执行的工艺决策}。}

人工编制的问题主要是重复劳动多、依赖个人经验、相似零件的结果不一致、文件维护困难。CAPP的价值是提高效率与一致性，集中保存工艺知识，并为CAD/CAM集成提供连接。

\begin{center}
\begin{tabularx}{\linewidth}{@{}p{.18\linewidth}XX@{}}
\toprule
接口 & 接收的信息 & 输出或反馈的信息\\
\midrule
CAD & 几何拓扑、材料、尺寸公差、粗糙度等 & 可制造性评价、设计修改建议\\
CAM & 工艺执行中的修改意见 & 加工顺序、设备工装、切削参数及相关刀位信息\\
生产管理 & 资源、生产状态及工艺修改意见 & 工时、成本、工艺文件和资源需求\\
质量系统 & 检测结果与质量反馈 & 检验工序、基准、检测要求与质量控制数据\\
\bottomrule
\end{tabularx}
\end{center}

\nt{几何模型能被两个软件打开，只说明具备某种数据交换能力。完整的CAD/CAPP/CAM集成还要求尺寸、公差、材料、特征与版本等信息的语义一致。单有三维外形不能唯一推出加工方案。}

\section{系统结构与零件信息输入}
\sourcepage{CAPP补充讲义}{5--10}

CAPP系统通常包含零件信息输入、工艺决策、工艺数据库与知识库、人机界面、文件管理与输出。控制模块协调各部分。工艺决策可分为三个层次：
\begin{enumerate}
\item \textbf{工艺路线：}确定加工阶段、工序顺序与所用设备，形成工艺过程卡。
\item \textbf{工序：}确定该次装夹中的定位夹紧、加工要求、工序尺寸及设备工装，形成工序卡。
\item \textbf{工步：}细化刀具、切削用量、走刀与质量要求，为刀位计算和NC程序准备数据。
\end{enumerate}

\subsection{输入信息应包括什么}
零件名称、图号与版本属于总体信息；点、边、面与连接关系属于几何拓扑信息；材料、公差、粗糙度、热处理等属于工艺相关要求。还须知道可用机床、刀具夹具、加工能力与企业规则，否则不能判断方案是否可行。

\begin{center}
\begin{tabularx}{\linewidth}{@{}p{.23\linewidth}XX@{}}
\toprule
描述或输入方法 & 基本思路 & 主要限制\\
\midrule
分类编码（GT） & 用代码描述零件类别和特征，按相似性分组 & 简单、适合派生式；长码易错，难完整表示复杂信息\\
语言及交互输入 & 用专门语言或逐项填写特征输入 & 需学习规则；重复录入耗时，可能丢失原设计信息\\
特征识别 & 从CAD几何拓扑中识别孔、槽、台阶等加工特征 & 相交特征有歧义；公差、热处理等不一定可由几何推知\\
基于特征建模 & 建模时建立可供下游读取的特征及属性 & 需要统一特征模型及可维护的特征库\\
产品数据交换 & 在共同产品信息模型中交换所需数据 & 实际能否传递语义取决于所用模型范围和接口实现\\
图纸扫描识别 & 从扫描图像恢复图元，再识别工艺信息 & 图像识别与特征识别均有误差和歧义\\
\bottomrule
\end{tabularx}
\end{center}

\ex{为什么识别出一个圆柱孔仍不够？}{相同直径和深度的孔，若一个只要求低精度，另一个要求严格的位置精度，候选工艺可能分别为钻削与钻后镗削。材料、孔轴线相对基准的位置要求和表面粗糙度也会改变选择。因此“识别形状”与“读懂完整制造要求”是不同任务。}

\section{派生式CAPP：从相似零件得到新工艺}
\sourcepage{CAPP补充讲义}{11--16}

\dfn{派生式CAPP}{基于成组技术（GT），利用“相似零件往往具有相似工艺”的规律，检索所属零件族的标准工艺，再按当前零件要求筛选、编辑和补充。也称检索式或变异式CAPP。}

\reviewfig[.78]{ch08-variant-workflow.pdf}{派生式CAPP的准备阶段与使用阶段（CAPP讲义，PDF第11页局部）。}{fig:capp-variant}

\subsection{准备与使用是两个不同阶段}
\textbf{准备阶段}选择分类编码体系，将零件分组，设计代表零件族特征的样件，建立标准工艺、工步代码及机床、刀具、余量、切削参数等文件。样件是零件族的抽象复合代表，不一定是某个实际生产件；标准工艺应综合覆盖族内所有零件的加工需求，再按当前零件筛选。

\textbf{使用阶段}输入新零件信息，定位零件族，检索标准工艺，再根据特征及技术要求保留或删除工步、修改参数并校核，最终输出当前零件的工艺文件。准备阶段建立“可供复用的知识”，使用阶段把它变成“具体零件的方案”。

\ex{原讲义工步代码筛选：11202}{代码前两位“11”表示精车外圆；第三位“2”表示检查零件分类码的第2位；末两位“0、2”给出允许范围$[0,2]$。若零件的JBM-1码为013124279，则第2位为1，满足$0\leq1\leq2$，保留该工步。这里判定的是\textbf{指定码位是否属于区间}，不能把工步代码11202当作与整串零件码比较的数值。若指定码位为3，则本筛选规则不保留该工步。}

派生式的优势是实现相对容易，能利用成熟经验，结果接近企业习惯；不足是依赖零件族、标准工艺和人工修改的质量。新零件超出既有族的范围，或企业资源发生变化时，直接套用旧路线可能失效。相似性是一种检索依据，不是制造可行性的证明。

\section{创成式CAPP：由知识与规则生成方案}
\sourcepage{CAPP补充讲义}{17--19}

\dfn{创成式CAPP}{根据零件特征、技术要求和制造资源，利用工艺决策逻辑进行搜索、计算与推理，生成工艺路线及工序内容。其主要依据是\textbf{可表达的制造知识和决策机制}，而非直接修改一份检索到的标准工艺。}

开发时首先限定适用零件类，分析基本表面和可选加工方法，收集加工能力、经济加工精度、设备工装数据，再归纳决策规则。软件阶段把这些规则实现为决策树、决策表、算法或推理系统。范围与知识不完整，自动生成结果就可能不完整；创成式并不等于对所有零件都能完全自动规划。

\subsection{工艺决策的三类问题}
\begin{itemize}
\item \textbf{选择：}在满足能力约束的候选项中选择毛坯、机床、刀具、夹具、加工方法和参数。
\item \textbf{规划：}安排加工阶段、工序和工步顺序，处理基准、装夹和先后依赖。
\item \textbf{计算：}确定余量、工序尺寸、切削用量、工时与成本，对候选方案进行比较。
\end{itemize}
工艺路线设计给出骨架，工序设计把骨架细化到可执行程度，两者需要相互校核。

\reviewfig[.98]{ch08-decision-tree.pdf}{装夹与孔加工方法的决策树示意（CAPP讲义，PDF第19页局部）；分支阈值仅表示课件示例。}{fig:capp-tree}

\subsection{决策树与决策表}
决策树用节点表示条件、分支表示判断结果、叶节点表示动作，便于直观看到判断路径。决策表按列组织规则，列出条件组合与对应动作，适合检查条件是否遗漏或互相冲突。二者可表达相同逻辑，但组织方式不同。

\ex{把孔加工选择写成可检查的规则}{以讲义示意为例：若要求低精度，可选钻削；若尺寸和形状精度要求高，还需区分位置精度要求，分别考虑钻后镗削或钻后铰削。真正应用时还须补入孔径、材料、余量、机床能力和前序基准。只写“高精度就铰孔”会遗漏位置精度约束，也没有覆盖所有输入情形。}

\nt{检查规则时至少问三件事：每个允许输入是否都有处理结果；同一输入是否触发矛盾动作；建议的设备和方法是否确实具备满足要求的能力。阈值来自特定工艺知识和资源，不能脱离企业条件照搬。}

\section{专家系统及AI辅助的定位}
\sourcepage{CAPP补充讲义}{20--21}

专家系统将\textbf{知识库}与\textbf{推理机}分离：知识库存储事实、规则和经验，推理机在控制策略下选择并应用知识。数据库保存材料、设备、参数等数据，解释模块说明某个结论是怎样得到的，知识获取模块负责补充和维护知识，人机界面连接用户。

一般算法式创成系统主要将决策逻辑写入算法和表；专家系统更强调知识的独立表示、搜索推理与解释。学习功能可以辅助知识获取，但不是任意规则系统天然具备的能力。

讲义提出的AI辅助场景可归为工艺路线推荐、参数与工时估算、设计可制造性反馈、变更影响分析、经验检索和文件校审。它们共同依赖完整的产品信息、企业资源数据和可追溯的工艺经验。推荐路线仍要检查加工能力、工序顺序、装夹和质量要求，不能以推荐分数替代工艺验证。

\Needspace{13cm}
\section{复习速查与自测}
\begin{center}
\begin{tabularx}{\linewidth}{@{}p{.2\linewidth}XX@{}}
\toprule
比较对象 & 派生式 & 创成式\\
\midrule
直接依据 & 零件族与标准工艺 & 特征、工艺知识与决策逻辑\\
主要操作 & 检索、筛选、编辑 & 搜索、推理、计算、规划\\
前期重点 & 分类编码、族划分、标准路线 & 加工能力与规则的建模\\
主要局限 & 超出族范围时复用能力弱 & 知识获取及规则覆盖困难\\
共同要求 & \multicolumn{2}{l}{完整输入、真实资源数据、结果校核与知识维护}\\
\bottomrule
\end{tabularx}
\end{center}

\begin{enumerate}
\item CAPP的输入是否只有CAD几何？\textit{答：还需材料、技术要求、资源、批量及工艺知识等。}
\item 为什么工艺设计是CAD与CAM的桥梁？\textit{答：把产品定义转化为制造方案和CAM需要的工艺条件。}
\item 样件是否必须实际存在？\textit{答：不必，它是零件族的抽象复合代表。}
\item 标准工艺检索后能否直接下发？\textit{答：应按当前零件和资源筛选、编辑并校核。}
\item 创成式与专家系统是否同义？\textit{答：专家系统是实现知识推理的一种方式，创成式也可用决策树、表和算法实现。}
\item 工步11202检查第几位？零件码第2位为3时如何？\textit{答：第2位；不在0至2内，本规则不保留该工步。}
\item 有三维模型但无法自动规划，可能缺什么？\textit{答：工艺语义、加工特征、资源能力或足够的决策知识。}
\end{enumerate}
